\documentclass[11pt]{article}
\usepackage{amsmath,amssymb,amsthm}
\usepackage{geometry}
\usepackage{hyperref}
\usepackage{microtype}
\usepackage{booktabs}

\geometry{margin=1.2in}

\newtheorem{theorem}{Theorem}[section]
\newtheorem{corollary}[theorem]{Corollary}
\newtheorem{lemma}[theorem]{Lemma}
\newtheorem{definition}[theorem]{Definition}
\newtheorem{remark}[theorem]{Remark}

\DeclareMathOperator{\logit}{logit}
\newcommand{\updateBelief}{\mathrm{updateBelief}}
\newcommand{\R}{\mathbb{R}}

\title{Finding the Bayes Net in the Transformer:\\
A Tutorial on Message Passing, Belief Propagation, and Transformer Circuits}
\author{Greg Coppola\\
\small coppola.ai \quad greg@coppola.ai}
\date{2026}

\begin{document}

\maketitle

\begin{abstract}
A growing body of work suggests that transformer computation can be
understood through the lens of probabilistic message passing. This paper
is a tutorial introduction to that perspective. Using the formally
verified result that sigmoid transformers can exactly realize belief
propagation~\cite{coppola2026transformers} as an organizing center, we
walk through the transformer architecture component by component:
residual streams as evolving belief states, attention as evidence
routing, feed-forward networks as belief updates, and layers as rounds
of message passing.

The goal is not to argue that all prior work is secretly ``the same
thing,'' nor to claim that every transformer literally instantiates an
exact Bayesian network in every detail. Rather, the goal is to provide a
common vocabulary for understanding why several previously separate
research programs --- graphical models, mechanistic interpretability,
Bayesian accounts of in-context learning, graph neural networks, and
distributed representations --- repeatedly converge on similar
structures.

Along the way, we develop the underlying log-odds algebra of Turing,
Good, and Pearl, discuss grounding and hallucination through the lens of
verification and ontology, and survey empirical findings from
interpretability research that naturally fit the message-passing view.
The companion paper~\cite{coppola2026transformers} develops the formal
constructive theorems; this paper is intended as a conceptual and
mechanistic guide to the broader picture.
\end{abstract}

\tableofcontents
\newpage

\part{The Mathematical Foundation}
\section{Introduction}
\label{sec:introduction}

A sigmoid transformer can exactly realize belief propagation
computations. This is not a loose analogy or a high-level functional
similarity: under the constructive mapping proved in the companion
paper~\cite{coppola2026transformers}, the transformer forward pass
implements the BP update equations directly. The present paper is a
tutorial guide to what that result means, how it connects to
transformer components, and what it implies for understanding trained
models.

The formal anchor is the companion theorem. This paper does not
rederive it. Instead, it builds the conceptual picture around it:
opening each transformer component, identifying its role in the
message-passing account, and connecting the formal result to a
converging body of empirical work in mechanistic interpretability,
in-context learning, and graphical model theory.

\subsection*{Scope of the Claims}

The paper operates at four distinct levels:

\begin{enumerate}
\item \textbf{Exact constructive results.} The companion theorem gives
an exact realization of BP in sigmoid transformers under a constructive
weight mapping~\cite{coppola2026transformers}.

\item \textbf{Empirical observations.} Mechanistic interpretability
results from trained transformers exhibit structures consistent with
the BP interpretation.

\item \textbf{Mechanistic interpretation.} Transformer components can
be understood through a routing-and-inference decomposition aligned
with BP primitives.

\item \textbf{Speculative extensions.} The discussions of grounding,
hallucination, and hybrid continuous-discrete inference are
interpretive extensions rather than established theorems.
\end{enumerate}

These claims have different evidentiary status and are marked
accordingly throughout the paper.

\subsection*{The Five Correspondences}

The tutorial is organized around five structural correspondences
between transformer components and the operations of belief
propagation:

\begin{enumerate}
\item \textbf{The residual stream is the belief state.} Each token's
$D_{\mathrm{model}}$-dimensional vector stores the current accumulated
state of inference at that node. Residual connections preserve and
update beliefs across layers rather than replacing them outright.

\item \textbf{Attention performs evidence gathering.} Attention heads
route information from neighboring tokens into the residual stream
before the update step executes. Multiple required evidence sources
become jointly available in a shared workspace.

\item \textbf{The FFN performs belief update.} Under the constructive
BP mapping, the feed-forward network computes
$\updateBelief(m_0, m_1) = \sigma(\logit(m_0) + \logit(m_1))$ from
the gathered evidence. This is the Bayesian update rule for
independent binary evidence sources --- exact with sigmoid activation,
approximate with ReLU or SwiGLU.

\item \textbf{Layers correspond to rounds of message passing.} One
layer is one gather-then-update cycle. Stacking layers corresponds
naturally to iterative rounds of inference.

\item \textbf{The scaling dimensions describe inference capacity.} The
architectural dimensions --- $W$, $L$, $H$, $D_{ff}$,
$D_{\mathrm{model}}$ --- determine the scale of the routing and
inference operations available to the network.
\end{enumerate}

\subsection*{The Running Example}

Throughout the paper we will repeatedly use a simple example involving
three propositions:
\[
\texttt{like(jack,jill)}, \quad \texttt{like(jill,jack)}, \quad
\texttt{date(jack,jill)}.
\]

The example is intentionally small. Its purpose is not realism but
clarity. We use it repeatedly to illustrate how attention gathers
evidence, how FFNs combine beliefs, how residual streams store
intermediate state, how layers propagate information over multiple
rounds, and what grounding means in a formally declared concept space.
By keeping the example fixed across sections, the paper can focus on
computational structure rather than constantly introducing new
notation.

\subsection*{The Underlying Algebra}

All five correspondences rest on the same mathematical spine: the
log-odds algebra of independent binary evidence, developed by Turing
and Good~\cite{good1950probability}, formalized by Pearl as the
sum-product algorithm~\cite{pearl1988probabilistic}, and realized
exactly by the constructive sigmoid mapping. Within this mapping,
sigmoid is the exact activation required to recover the log-odds
update equations. Part~I develops this foundation before turning to
the architectural correspondences.

\subsection*{Grounded and Ungrounded Inference}

The distinction between grounded and ungrounded transformers is not
primarily architectural. In both cases, the network performs inference
over an implicit factor graph. The difference is whether the graph has
a declared ontology, finite concept space, and verifier that permits
formal correctness guarantees.

For grounded transformers operating on tree-structured knowledge
bases, the constructive BP results provide exact posterior guarantees.
Standard language models do not provide guarantees of this kind
because their internal representations are not tied to a formally
declared ontology. Part~IV develops this distinction and its
implications for hallucination and reliability.

\subsection*{Convergent Evidence}

Five independent research programs have converged on the same general
picture from different starting points: the transformer forward pass
implements a form of distributed probabilistic inference. Mechanistic
interpretability studies identify routing heads and structured FFN
updates. Belief-state geometry work finds posterior-like organization
in residual streams. In-context learning work models transformers as
Bayesian learners. Graph neural network research established the
message-passing connection. Constructive weight-learning experiments
show that optimization reliably discovers BP-like solutions. Part~V
surveys this convergence.

\subsection*{How to Read This Paper}

Part~I establishes the mathematical foundation. Readers already
familiar with log-odds algebra and belief propagation can skim
portions of it. Part~II develops the core architectural
correspondences component by component. Part~III examines
superposition, the three softmaxes, and the relationship between
trained weights and implicit factor graphs. Part~IV addresses
grounding and hallucination. Part~V surveys convergent evidence and
related work.

The formal constructive results are developed
in~\cite{coppola2026transformers}. This paper is a tutorial guide to
the conceptual landscape those results open up.
\section{The Log-Odds Algebra}
\label{sec:log-odds}

The mathematical spine of the entire account is a single algebraic structure: the group of beliefs on $(0,1)$ under log-odds addition. This algebra was discovered three times independently, in three different contexts, by three different research traditions. Each time, the same formula emerged.

\subsection*{First Discovery: Turing and Good at Bletchley Park}

During World War II, Alan Turing and I.J.\ Good needed a practical method for combining independent pieces of evidence when breaking Enigma. The method they developed was log-odds addition~\cite{good1950probability}.

Given a binary hypothesis $H$ and independent evidence sources $e_0, e_1$, the weight of evidence adds:
\[
\logit(P(H \mid e_0, e_1)) = \logit(P(H)) + W(H : e_0) + W(H : e_1).
\]
Starting from a uniform prior $P(H) = 0.5$, so $\logit(P(H)) = 0$, and setting $m_i = P(H \mid e_i)$:
\[
\logit(P(H \mid e_0, e_1)) = \logit(m_0) + \logit(m_1).
\]
Converting back via sigmoid: $P(H \mid e_0, e_1) = \sigma(\logit(m_0) + \logit(m_1)) = \updateBelief(m_0, m_1)$.

The updateBelief function is Turing and Good's weight-of-evidence combination, applied to two binary evidence sources. It is not an approximation. For independent evidence sources, it is exact.

\subsection*{Second Discovery: Pearl's Sum-Product Algorithm}

Pearl~\cite{pearl1988probabilistic} derived the belief propagation update rule for binary variable nodes from the general sum-product equations applied to pairwise factors with independent messages. The result is:
\[
b_{\mathrm{new}}(v) = \frac{m_0 m_1}{m_0 m_1 + (1-m_0)(1-m_1)} = \sigma(\logit(m_0) + \logit(m_1)).
\]
Pearl derived this from the structure of graphical models. He did not frame it as log-odds addition or connect it to Turing and Good's weight-of-evidence tradition. The formula is the same. The framing is new.

Pearl's deeper contribution was showing that this update rule, applied
locally at each node using only messages from immediate neighbors, is
sufficient to compute exact marginal posteriors on tree-structured
graphs~\cite{pearl1988probabilistic}. Global inference emerges from
local computation. No central coordinator is needed: each node updates
its belief from its neighbors' current beliefs, and iterating this
process to convergence yields the exact answer. This is the
sum-product algorithm, and it is what makes the connection to
transformers non-trivial. A transformer is also a distributed
computation in which each position updates from its neighbors across
layers, with no global state outside the residual stream. The
architectural parallel is not incidental --- it is the reason the same
formula appears in both systems.

\subsection*{Third Discovery: The Sigmoid Transformer}

The sigmoid FFN computing $\sigma(w_0 \cdot \logit(m_0) + w_1 \cdot \logit(m_1) + b)$ is performing weighted log-odds addition and converting back to probability space in a single operation. With $w_0 = w_1 = 1$ and $b = 0$, this is $\updateBelief$ exactly. The sigmoid activation is not a design choice motivated by gradient flow or output normalization. It is the exact function required to implement the Turing-Good-Pearl algebra.

\subsection*{The Algebraic Structure}

Log-odds addition defines a commutative group on $(0,1)$ via the isomorphism $\logit : (0,1) \to \mathbb{R}$:
\[
m_0 \oplus m_1 = \sigma(\logit(m_0) + \logit(m_1)).
\]

\begin{itemize}
\item \textbf{Commutativity:} $m_0 \oplus m_1 = m_1 \oplus m_0$.
\item \textbf{Associativity:} $(m_0 \oplus m_1) \oplus m_2 = m_0 \oplus (m_1 \oplus m_2)$.
\item \textbf{Identity:} $m \oplus 0.5 = m$, since $\logit(0.5) = 0$.
\item \textbf{Inverse:} $m \oplus (1-m) = 0.5$, since $\logit(m) + \logit(1-m) = 0$.
\end{itemize}

The identity element is $0.5$ --- maximum uncertainty, no information. Combining any belief with a $0.5$ belief leaves it unchanged. This is the formal statement that a neutral prior contributes nothing.

\subsection*{Relation to Classical Boolean Logic}

This algebra is the probabilistic generalization of classical Boolean AND and OR. In the limit as beliefs approach 0 and 1:
\begin{itemize}
\item $\logit(1) = +\infty$ (certain true), $\logit(0) = -\infty$ (certain false).
\item Two large positive numbers sum to a larger positive number: high confidence in both inputs yields high confidence in their conjunction.
\item Equal and opposite evidence cancels: $\logit(p) + \logit(1-p) = 0$, yielding $0.5$ --- maximum uncertainty.
\end{itemize}

Classical Boolean logic is the limiting case as beliefs become certain. The log-odds algebra does not have an annihilator (FALSE AND TRUE $\neq$ FALSE in the probabilistic case); instead it has cancellation: strong evidence in opposite directions yields maximum uncertainty. This is the correct behavior for uncertain reasoning.

\subsection*{Why This Matters}

The three independent rediscoveries of the same formula are not a coincidence. The log-odds algebra is the unique correct way to combine independent binary evidence under the axioms of probability theory. Any system that does probabilistic reasoning over binary propositions with independent evidence sources will converge on this algebra. The transformer converged on it because it was trained to do probabilistic reasoning. Gradient descent did not discover something new. It rediscovered the structure that the analysis of reasoning already required.
\section{Two Theorems, One Weight Construction}
\label{sec:turing-complete}

Before developing the five correspondences in detail, it is worth pausing on a
structural fact that makes the BP interpretation more than a clever observation.
The transformer's routing mechanism is so clean that the \emph{same two weight
families} appear in two completely independent proofs: the Turing completeness
proof and the BP proof. This is not a coincidence. It reveals that the
projectDim/crossProject construction is the natural primitive for what attention
heads actually do.

\subsection*{Turing Completeness}

P\'{e}rez et al.~\cite{perez2021turing} proved that transformers are Turing
complete under floating-point arithmetic. The companion
paper~\cite{coppola2026transformers} gives a constructive Lean-verified proof
via Boolean circuit simulation:

\begin{theorem}[Transformer is Turing Complete]
For any Turing machine $M$, there exist transformer weights $W$ such that for
all tapes $\tau$ and step counts $t$:
\[
\mathrm{transformerAgent}(W, \mathrm{encode}(M, \tau), t)
= \mathrm{encode}(M.\mathrm{run}(\tau, t)).
\]
Formally verified in Lean~4 against standard mathematical axioms.
\end{theorem}

The proof is constructive. Any Turing machine step decomposes into a Boolean
circuit. A transformer simulates any Boolean circuit in a single FFN forward
pass via threshold gates. The key insight is the role of each component:
\emph{attention is lookup} (find the tape cell at the current head position),
\emph{FFN is gating} (apply the transition function), and \emph{the agent loop
is iteration} (step the machine).

\subsection*{The Same Two Weight Families}

The Turing completeness proof and the BP proof use the same two sparse matrix
families for Q/K and V weights:

\begin{itemize}
\item $\mathrm{projectDim}(d)$: a rank-1 diagonal matrix extracting dimension
$d$. Used as $W_Q$ and $W_K$. The attention score $e_j[d] \cdot e_k[d]$ peaks
when token $k$'s stored index matches token $j$'s query value --- exact index
matching. In the TM proof: finds the tape cell at the head position. In the BP
proof: finds the neighbor node whose belief to gather.

\item $\mathrm{crossProject}(s, d)$: a rank-1 off-diagonal matrix reading
source dimension $s$ and writing to destination dimension $d$. Used as $W_V$.
In the TM proof: routes the tape symbol to the transition function. In the BP
proof: routes the neighbor's belief into the residual stream scratch slot.
\end{itemize}

Both proofs also reuse two results from the Turing completeness Lean
development: \texttt{posEncDot\_distinct} (the $Q \cdot K$ score is strictly
maximized at the correct target) and \texttt{softmax\_concentrates} (at
sufficient temperature, softmax converges to a point mass at the argmax).

\subsection*{What This Means}

The shared weight construction is significant for two reasons.

First, it means the routing mechanism is \emph{canonical}. There is one natural
way to make attention heads perform index-based lookup in a transformer, and
both the Turing and BP interpretations use it. This is not a special-case
construction tailored to BP --- it is the architecture's natural primitive.

Second, it separates routing from inference. The projectDim/crossProject
construction handles routing (which token to look at) identically in both
proofs. What differs is what happens \emph{after} the lookup: the TM proof
uses the FFN as a Boolean gate; the BP proof uses it as a probabilistic
update. Same routing, different inference. The transformer is a general-purpose
routing-plus-inference architecture, and both Turing completeness and BP are
natural instances of it.

The Turing completeness result says: the transformer \emph{can} compute
anything. The BP result says: it \emph{does} compute something specific ---
belief propagation --- with these specific weights. Both are needed for a
complete picture of what the architecture is capable of and what it actually
does.

\subsection*{Empirical Confirmation}

The \texttt{learner} repository confirms that gradient descent finds
TM-simulation weights from scratch. Five structurally distinct Turing machines
--- binary incrementer, decrementer, bitwise complement, left shift, right
shift --- were trained using a standard transformer encoder (2~layers, 2~heads,
$d_{\mathrm{model}} = 32$, $\sim$4K parameters) from random initialization.
All five reached 100\% validation accuracy in 4 epochs with identical
hyperparameters and no per-machine tuning~\cite{coppola2026transformers}.

The formal construction is not merely theoretically possible --- it is what
gradient descent finds given a clean training signal. This parallels the
\texttt{bayes-learner} result for BP (val MAE 0.000752), which we develop in
Section~\ref{sec:ffn-or}. Together they establish a consistent pattern: the
formal weight constructions are what the architecture naturally learns.

\part{The Five Correspondences}
\section{The Residual Stream as the Belief State}
\label{sec:residual-stream}

The starting point for seeing the Bayes net in the transformer is the residual
stream. Everything else follows from getting this right.

In the standard transformer encoder, each token position carries a vector of
dimension $D_{\mathrm{model}}$ that persists across all layers. Each layer reads
from this vector and writes its output back into it via a residual connection.
Nothing is discarded; the stream accumulates. Elhage et al.~\cite{elhage2021mathematical}
name this the \emph{residual stream} and frame it as a shared workspace: attention
heads and feed-forward layers are read-write operations on a common scratchpad.

In the Bayes net view, the residual stream at token position $t$ is the
\emph{belief state} of factor graph node $t$. Specifically, the formal token
encoding from~\cite{coppola2026transformers} assigns clean semantic roles to each
dimension:

\begin{center}
\begin{tabular}{cl}
\toprule
Dimensions & Content \\
\midrule
0 & own belief (initialized to $0.5$) \\
1--4 & factor table entries $[f_{00}, f_{01}, f_{10}, f_{11}]$ \\
5 & node type (0 = variable, 1 = factor) \\
6 & own index $/ (n-1)$ \\
7 & neighbor index $/ (n-1)$ \\
\bottomrule
\end{tabular}
\end{center}

Dimensions 5--7 form the \emph{routing key}: the complete inferential identity of
the token. Dimensions 0--4 are continuous parameters: they supply the magnitudes
the inference computes over, but they do not determine the structure of the
computation. This split is not a convenience of presentation --- it is a theorem.
Two tokens with the same routing key receive identical attention patterns
regardless of their continuous parameters~\cite{coppola2026transformers}.

\subsection*{Routing Key = Concept. Continuous Dimensions = Magnitude.}

This is the right level of abstraction for thinking about what the residual stream
represents. The routing key identifies \emph{which} concept a token corresponds to.
The continuous dimensions encode \emph{how strongly} the evidence supports that
concept. A production transformer with $D_{\mathrm{model}} = 4096$ is not tracking
4096 independent propositions --- it is tracking a much smaller number of concepts
(bounded by the routing key space) while using the remaining dimensions to encode
magnitudes and handle superposition.

The residual connection is what makes this interpretation coherent across layers.
Belief propagation is an iterative algorithm: each round refines the belief at every
node by incorporating messages from neighbors. The residual connection ensures that
each layer's update is added to the accumulated belief, not computed from scratch.
Layer $l$ does not replace the belief; it refines it.

\subsection*{What This Looks Like in Trained Models}

The mechanistic interpretability literature has independently converged on a
compatible picture. Elhage et al.~\cite{elhage2022superposition} show that the
residual stream represents far more features than $D_{\mathrm{model}}$ dimensions
via superposition: features are encoded as nearly-orthogonal directions in a
high-dimensional space, with interference managed by sparsity. In the Bayes net
view, this means the residual stream is simultaneously tracking many factor graph
nodes, each encoded as a direction rather than a dedicated dimension.

Sparse autoencoder (SAE) features~\cite{templeton2024monosemanticity} are the
natural candidates for individual factor graph nodes. A monosemantic SAE feature ---
one that activates for a single coherent concept --- corresponds to a single routing
key in the formal construction. The feature direction in residual stream space is
where the belief value for that concept lives.

The broader implication: the residual stream is not a flat representation. It has
structure --- routing structure, imposed by the attention weights, that determines
which features interact with which. That structure is the implicit factor graph.
\section{Attention as Evidence Gathering}
\label{sec:attention-and}

Within the BP interpretation, the attention mechanism performs the
gather phase of message passing: it routes neighboring beliefs into the
residual stream before the update step executes. This corresponds to the
evidence-gathering phase of Pearl's sum-product
algorithm~\cite{pearl1988probabilistic}.

The conjunctive structure is architectural rather than localized inside
any individual attention head. Consider a single transformer layer. Each
attention head scans token positions, concentrates on a relevant
neighbor, and copies information from that neighbor into a designated
region of the current token's residual stream. Only after all heads have
finished does the feed-forward network execute.

The conjunction is therefore not inside any one head. It emerges at the
level of the residual stream, where multiple required evidence sources
become simultaneously available in a shared workspace before the update
step runs. The FFN has no mechanism to execute on partial information:
it reads the entire residual stream state. In this sense, the
architecture enforces a gather-before-update semantics analogous to
conjunctive evidence integration.

\subsection*{The Formal Weight Construction}

The constructive mapping from
\cite{coppola2026transformers} realizes the gather step exactly. Each
attention head uses two sparse weight families:

\begin{itemize}
\item $\mathrm{projectDim}(d)$: a diagonal projection extracting
dimension $d$. Used as $W_Q$ and $W_K$. The resulting $Q \cdot K$ dot
product peaks when token $k$'s stored index matches token $j$'s query
value, implementing exact index matching.

\item $\mathrm{crossProject}(s, d)$: an off-diagonal projection reading
source dimension $s$ and writing to destination dimension $d$. Used as
$W_V$. Routes neighbor $k$'s belief from one residual-stream dimension
into another.
\end{itemize}

At sufficiently high temperature, the softmax concentrates almost
entirely on the matching neighbor. In the constructive setting, sharp
attention behaves like a routing or lookup operation rather than an
inference computation. The probabilistic update occurs in the FFN stage.

\subsection*{A Concrete Example}

Consider three propositions:
\texttt{like(jack,jill)} with belief $m_0 = 0.8$,
\texttt{like(jill,jack)} with belief $m_1 = 0.4$, and
\texttt{date(jack,jill)} whose belief is to be updated.

Head~0 attends to \texttt{like(jack,jill)} and writes $0.8$ into one
region of the \texttt{date} token's residual stream. Head~1 attends to
\texttt{like(jill,jack)} and writes $0.4$ into another region. By the
time the FFN executes, both evidence sources are simultaneously
available. The gather phase is complete before the update step begins.

\subsection*{Evidence from Trained Models}

Mechanistic interpretability studies have identified attention patterns
consistent with the routing behavior predicted by the constructive
mapping. Wang et al.~\cite{wang2022ioi} identify ``name mover heads'' in
GPT-2 small that copy token-specific information to target positions,
behavior closely resembling crossProject-style routing. Their Q/K
circuits also exhibit approximately rank-1 structure, qualitatively
consistent with projectDim-style matching.

Empirical classification studies further identify several recurring head
types:

\begin{itemize}
\item \textbf{Semantic routing heads}: attend to semantically related
tokens determined by content rather than position. These resemble the
gather step of the constructive BP mapping.

\item \textbf{Hub heads}: attend consistently to fixed positions,
especially the BOS token. These appear to distribute or collect global
context information through a shared communication channel.
\end{itemize}

These observations do not by themselves prove the BP interpretation, but
they are structurally consistent with a routing-plus-update view of
transformer computation.

\subsection*{Scaling}

For Llama~3 405B, there are $L \times H = 126 \times 128 = 16{,}128$
distinct attention-head routing patterns. The number of routing
operations per forward pass scales as
$L \times H \times W = 132\mathrm{M}$ at sequence length
$W = 8192$. Within the BP interpretation, these routing operations form
the communication structure of the implicit factor graph.
\section{The FFN as Belief Update}
\label{sec:ffn-or}

Once attention has assembled the relevant evidence into the residual
stream, the feed-forward network performs the update step. In the exact
constructive setting, this update is the binary belief-propagation rule
for combining independent evidence sources.

The central formula is:
\[
\updateBelief(m_0, m_1)
  = \frac{m_0 m_1}{m_0 m_1 + (1-m_0)(1-m_1)}
  = \sigma\!\left(\logit(m_0) + \logit(m_1)\right).
\]
With sigmoid activation and BP weights ($w_0 = w_1 = 1$, $b = 0$), the
FFN computes this exactly from the gathered beliefs in dimensions~4
and~5, writing the result to dimension~0. This is the Bayesian update
rule for two independent binary evidence sources in the constructive
model~\cite{coppola2026transformers}.

\subsection*{The Boolean Structure Inside the Update}

The formula contains the same structure that appears in classical
Boolean reasoning, but in probabilistic form. The term $m_0 m_1$ is the
joint probability that both independent inputs are true. In log-odds
space, this multiplication becomes addition:
\[
\logit(\updateBelief(m_0, m_1)) = \logit(m_0) + \logit(m_1).
\]

Independent evidence adds. The sigmoid converts the log-odds sum back to
probability space. The FFN is therefore a weight-of-evidence combiner in
the sense of Turing and Good~\cite{good1950probability}, when its inputs
and weights instantiate the constructive BP representation.

This is why the OR-gate language is useful: the FFN combines gathered
evidence into an updated conclusion. But the exact theorem is more
specific than the metaphor. The FFN implements the BP update exactly in
the sigmoid construction; production FFNs approximate or generalize this
role rather than literally instantiating the same binary rule in every
unit.

\subsection*{The Key-Value Memory Interpretation}

Geva et al.~\cite{geva2021ffn} show that FFN layers can be viewed as
key-value memories: the rows of $W_1$ are keys that detect input
patterns, and the columns of $W_2$ are values written to the residual
stream when those patterns activate. Each hidden unit detects a pattern
and contributes an update.

In the BP interpretation, these key-value memories resemble learned
factor potentials. The key determines when an update is relevant; the
value determines what direction is written back into the residual
stream. The full FFN layer integrates many such updates into the token's
next residual-stream state.

ROME~\cite{meng2022rome} provides suggestive evidence for this
interpretation: editing localized FFN parameters can change factual
associations in targeted ways. This does not by itself prove that FFN
weights are literal conditional-probability tables, but it is consistent
with the view that FFNs store semantically structured update directions.

\subsection*{Sigmoid vs.\ ReLU}

The formal proof requires sigmoid activation because $\updateBelief$ is
exactly $\sigma(\logit(m_0) + \logit(m_1))$: sigmoid is the inverse of
logit, and the composition gives the desired update. ReLU, GeLU, and
SwiGLU do not implement this binary log-odds algebra exactly.

Nevertheless, these activations may still support approximate or
generalized message-passing behavior. They preserve smooth, structured
transformations of residual-stream features, and empirical experiments
show that gradient descent can find approximately correct BP-like
solutions even with ReLU activations. In the \texttt{bayes-learner}
experiment, for example, the learned model reached validation MAE
$0.000752$~\cite{coppola2026transformers}.

Thus sigmoid gives the exact algebraic identity. Modern activations give
a practical approximation or extension, not the same formal guarantee.

\subsection*{Scaling}

For Llama~3 405B, there are
$L \times D_{ff} = 126 \times 53{,}248 \approx 6.7\mathrm{M}$
FFN hidden units across layers. The number of FFN evaluations per
forward pass is
$L \times D_{ff} \times W \approx 55\mathrm{B}$ at sequence length
$W = 8192$.

Within the BP interpretation, this large FFN budget reflects the cost of
implementing rich update rules over many superposed features. The
architecture devotes far more capacity to updating internal state than
to routing attention alone, which is consistent with the view that
inference is the dominant computational burden.
\section{Layers Are Rounds of Belief Propagation}
\label{sec:layers-rounds}

One transformer layer is one round of belief propagation. $L$ layers are $L$ rounds.
This is the content of Theorem~1.1 in~\cite{coppola2026transformers}, and it is the
claim that ties the previous three sections together.

A single round of BP proceeds in two steps. In the gather step, each variable node
copies its neighbors' current beliefs into scratch slots. In the update step, each
variable node computes a new belief from the gathered evidence. One transformer
layer does exactly this: attention is the gather step, the FFN is the update step,
and the strict alternation is architecturally enforced. The depth of the network is
not an engineering hyperparameter --- it is determined by the number of reasoning
hops required.

\subsection*{The Peirce Lower Bound}

No local algorithm can solve $N$-hop reasoning in fewer than $N$ rounds. This is
the Peirce lower bound, and it applies directly to belief propagation: on a tree of
diameter $d$, BP requires exactly $d$ rounds to propagate evidence from one end to
the other~\cite{pearl1988probabilistic}. A transformer with $L$ layers can perform
at most $L$ hops of reasoning. More layers are required for longer reasoning chains,
not for more parameters per se.

This reframes the question of why large models need depth. It is not that deeper
models have more capacity in a vague sense. It is that deeper models can perform
more rounds of inference --- they can follow longer chains of implication. A model
that needs to conclude $A \Rightarrow B \Rightarrow C \Rightarrow D$ needs at least
three layers, regardless of how wide each layer is.

\subsection*{Three-Phase Layer Structure}

Experiments classifying attention heads in GPT-2 small by behavior reveal a
three-phase structure across layers:

\begin{itemize}
\item \textbf{Layers 0--2 (continuous-heavy):} local feature extraction. Attention
heads attend to nearby tokens; the residual stream accumulates basic syntactic and
semantic features. This is the first round of BP: each node gathers information
from its immediate neighbors.

\item \textbf{Layer 3 (transition):} a shift in the dominant computation. Hub heads
begin to dominate. The BOS token starts accumulating a global context representation.

\item \textbf{Layers 4--11 (hub-dominated):} global context assembly. Hub heads
distribute the BOS representation to every token. This is iterative refinement of
a global prior --- each round incorporates the accumulated evidence from all previous
rounds.
\end{itemize}

The semantic-AND heads do not scale in number: 9 in GPT-2 small, 9 in GPT-2 medium.
What scales is communication infrastructure (hub heads), not logical reasoning
capacity per layer. This is consistent with the BP view: the number of reasoning
steps is set by depth, not width.

\subsection*{Residual Connections as Belief Accumulation}

The residual connection has a natural interpretation in this account. Each layer
adds its update to the accumulated belief, rather than replacing it. This is exactly
the structure of iterative BP: each round refines the beliefs from the previous
round. The residual stream after $l$ layers encodes the belief state after $l$ rounds
of message passing.

Without residual connections, each layer would compute its output from scratch and
the iterative refinement interpretation would break down. With residual connections,
the stream has a well-defined semantic content at every depth: it is the belief state
of the implicit factor graph after that many rounds of BP.

\subsection*{Convergence}

On tree-structured factor graphs, BP converges in exactly $\mathrm{diameter}(T)$
rounds and the resulting beliefs are the exact marginal posteriors. This is the
no-hallucination corollary of~\cite{coppola2026transformers}: a transformer with BP
weights, run for sufficiently many layers over a grounded tree-structured knowledge
base, computes exact Bayesian posteriors at every node.

On loopy graphs, BP may not converge, but in practice it does on the graph structures
that arise in typical knowledge bases. The \texttt{loopy} repository reports 100\%
convergence across 500 trials on five loopy graph structures, with mean KL divergence
below $0.0002$~\cite{coppola2026transformers}. The theoretical gap between trees and
loopy graphs is smaller in practice than in the worst case.
\section{The Scaling Numbers as a Bayes Net Inventory}
\label{sec:scaling-numbers}

The five architectural dimensions of a transformer --- $W$ (sequence length),
$D_{\mathrm{model}}$ (embedding dimension), $H$ (attention heads), $D_{ff}$ (FFN
hidden dimension), $L$ (layers) --- are not arbitrary engineering choices. In the
Bayes net view, each dimension has a precise interpretation, and together they
constitute an inventory of the implicit factor graph.

\begin{center}
\begin{tabular}{lll}
\toprule
Dimension & BP interpretation & Llama~3 405B \\
\midrule
$W$ & nodes in the factor graph (sequence positions) & 8{,}192 \\
$L$ & rounds of belief propagation & 126 \\
$H$ & AND patterns per layer & 128 \\
$D_{ff}$ & OR patterns per layer & 53{,}248 \\
$L \times H$ & distinct AND weight patterns & 16{,}128 \\
$L \times D_{ff}$ & distinct OR weight patterns & 6.7M \\
$L \times H \times W$ & AND evaluations per forward pass & 132M \\
$L \times D_{ff} \times W$ & OR evaluations per forward pass & 55B \\
\bottomrule
\end{tabular}
\end{center}

The ratio $D_{ff} / H \approx 416$ for Llama~3 405B. The architecture devotes
$416\times$ more capacity to inference (OR gates) than to routing (AND gates).
This is the right ratio for a system whose bottleneck is combining evidence, not
finding it.

\subsection*{What Scaling Adds}

When a model scales from 8B to 405B parameters, what changes in Bayes net terms?

\begin{itemize}
\item \textbf{More rounds of BP} ($L$ increases): deeper reasoning chains become
possible. A 126-layer model can follow implications 126 hops deep.

\item \textbf{Richer OR patterns} ($D_{ff}$ increases): more conditional probability
table entries per layer. The model can represent more fine-grained conditional
relationships.

\item \textbf{More nodes} ($W$ increases): longer sequences mean larger factor
graphs. The model reasons over more propositions simultaneously.

\item \textbf{AND patterns stay small} ($H$ grows slowly): 128 heads in 405B vs.\
16 in the smallest models. The number of distinct routing operations is nearly
constant. What scales is inference capacity, not routing capacity.
\end{itemize}

The implication is sharp: scaling adds inference depth and richness, not
fundamentally new reasoning structure. The AND/OR architecture is fixed. What
changes is how many rounds it runs and how many patterns it can match.

\subsection*{The 16,128 Number}

The 16,128 distinct AND weight patterns for Llama~3 405B is a striking number
because it is so small. Every token position in every sequence uses one of these
16,128 routing patterns --- the same patterns, regardless of what the token is.
The routing key (which neighbor to attend to) is determined by the weight matrices,
which are shared across positions. What varies across positions is the continuous
content of the residual stream, not the routing structure.

This is the formal statement of weight sharing as universal rule implementation
(Section~4.3 of~\cite{coppola2026transformers}): the same computation runs for every
instantiation of a concept, and the particular input supplies the magnitudes. The
16,128 patterns are the universal rules. The residual stream values are the
particular instances.

\part{Going Deeper}
\section{The Residual Stream vs.\ Propositions: Superposition and the $D_{ff}/D$ Ratio}
\label{sec:superposition}

The formal BP construction uses $D_{\mathrm{model}} = 8$ with one dedicated dimension per
proposition and one OR gate per belief update. A production transformer has $D_{\mathrm{model}}
= 4096$ and $D_{ff} = 14{,}336$: approximately 3.5 OR gates per residual stream dimension per
layer. This raises two related questions. Why does the ratio $D_{ff}/D$ exceed 1 at all? And
does $D_{ff} \approx 14{,}336$ mean there are 14,336 factor nodes per layer? This section
resolves both.

\subsection*{Hidden Units Are Not Factor Nodes}

In the minimal construction, one FFN hidden unit implements one factor node --- one
conditional relationship between two propositions. Hidden units and factor nodes are in
one-to-one correspondence.

In a production transformer, this correspondence breaks down in two ways:

\begin{enumerate}
\item \textbf{Multiple units per factor.} A single conditional relationship may require
multiple hidden units to represent accurately with non-sigmoid activations. ReLU and SwiGLU
require more units to approximate the exact log-odds combination that sigmoid implements in
one unit. The factor node is still the unit of semantics; the hidden units are the
implementation substrate.

\item \textbf{Superposed factors.} Because the residual stream represents many concepts via
superposition, a single hidden unit may participate in multiple factor relationships
simultaneously. The unit fires for a direction in residual stream space that is a
superposition of several feature directions, and the contribution it writes is a superposition
of several belief updates.
\end{enumerate}

The right level of abstraction is not hidden units but \emph{SAE features}. Each monosemantic
SAE feature corresponds to one factor node in the implicit factor graph. The number of
semantically distinct factor nodes per layer is the number of active SAE features --- typically
much smaller than $D_{ff}$.

$D_{ff}$ is therefore the \emph{implementation budget} for the OR gates, not the factor node
count. It sets an upper bound on representable conditional relationships per layer, but the
actual semantic content is determined by the SAE feature structure.

\subsection*{Three Reasons the Budget Exceeds One}

Three independent factors drive $D_{ff}/D$ above the minimal value of 1:

\textbf{Superposition.} Elhage et al.~\cite{elhage2022superposition} show that residual
streams represent far more features than $D_{\mathrm{model}}$ dimensions via superposition:
features are encoded as nearly-orthogonal directions in a high-dimensional space, with
interference managed by sparsity. A model with $D = 4096$ may track tens of thousands of
distinct concepts simultaneously. If the residual stream tracks $K$ active concepts per token
and each needs an independent belief update per layer, $D_{ff}$ must be at least $K$. Since
$K \gg D$, the OR budget must exceed the residual stream dimension by the superposition
density --- empirically, roughly $3$--$4\times$.

\textbf{$k$-Ary neighbors.} The formal construction handles pairwise factors: each OR gate
takes exactly two inputs. Real knowledge graph nodes have more than two neighbors. The scaling
theorem of~\cite{coppola2026transformers} handles this: any $k$-ary factor graph binarizes
exactly via a chain of $\lceil \log_2 k \rceil$ pairwise updates. A node with 8 neighbors
requires 3 chained pairwise updates rather than 1, contributing a further multiplicative
overhead to the required $D_{ff}$.

\textbf{Activation precision.} Sigmoid implements the log-odds algebra exactly. ReLU and
SwiGLU are compatible but not exact: they preserve non-negativity and smooth aggregation but
require more hidden units to achieve the same approximation quality. The bayes-learner
experiment confirms that gradient descent finds approximately correct BP solutions even with
non-sigmoid activations (val MAE 0.000752), but at the cost of additional hidden units per
factor.

Together these three factors explain why the minimal $D_{ff}/D = 1$ becomes $3$--$4$ in
production models. The ratio is an architectural constant encoding the combined implementation
overhead of superposition density, neighbor fanout, and activation imprecision.

\subsection*{The Dimension Budget and Factor Graph Scale}

A cleaner way to think about the implicit factor graph uses the dimension budget directly.
The residual stream has $D_{\mathrm{model}}$ dimensions. Each active concept occupies a
direction in this space. Empirically, roughly $10 \times D_{\mathrm{model}}$ features can be
superposed with manageable interference~\cite{elhage2022superposition}.

For Llama~3 405B with $D_{\mathrm{model}} = 16{,}384$, this gives roughly 160,000
simultaneously active features per token position. Each feature is a factor graph node. The
factor graph has on the order of $W \times 160{,}000 \approx 1.3\mathrm{B}$ nodes for a
sequence of length 8192. This is the scale of the implicit Bayesian network that one forward
pass of Llama~3 405B is performing inference over.

\subsection*{The OR/AND Asymmetry}

With this clarification, the $D_{ff}/H \approx 416$ ratio for Llama~3 405B has a clean
interpretation. It is not that there are 416 OR gates for every AND gate. It is that the OR
implementation requires 416 units of hidden dimension for every unit of attention head
dimension. The factor graph itself may be much sparser --- each node may have only 2--5
neighbors in practice --- but implementing those sparse connections accurately in the
superposed representation requires the $416\times$ overhead.

The implication for interpretability: finding the factor graph requires working at the SAE
feature level, not the hidden unit level. The factor potentials are encoded in the relationship
between input SAE features (what the $W_1$ rows respond to) and output SAE features (what the
$W_2$ columns write). The $D_{ff}$ hidden units are the implementation substrate; the SAE
features are the semantics.
\section{The Three Softmaxes}
\label{sec:three-softmaxes}

There are three places in the transformer where softmax or sigmoid appears. They do
three completely different jobs. Conflating them is the single most common source of
confusion about what the transformer is computing.

\begin{center}
\begin{tabular}{llll}
\toprule
& Job & Input & Semantics \\
\midrule
Softmax 1 & Routing & $Q \cdot K$ scores & Differentiable argmax \\
Softmax 2 & Inference & Neighbor beliefs & Bayesian posterior \\
Softmax 3 & Generation & Final hidden state & Token distribution \\
\bottomrule
\end{tabular}
\end{center}

\subsection*{Softmax 1: Attention (Routing)}

The first softmax appears inside every attention head, applied to the query-key dot
products. Its job is routing: which token should I look at? This is a differentiable
argmax --- at high temperature it concentrates all weight on the token with the
highest $Q \cdot K$ score.

In the BP construction, the attention softmax fetches a neighbor's belief by index
matching. The query encodes ``I am looking for the token at index $i$''; the key
encodes ``I am the token at index $i$''; the dot product peaks at the correct match;
softmax concentrates the weight there. The result is a lookup table.

This softmax is \emph{not} doing inference. It is not computing a probability over
hypotheses. It is selecting which piece of evidence to retrieve. Calling it a
``probability distribution over tokens'' is technically correct but misleading ---
the weights are routing coefficients, not posterior beliefs.

\subsection*{Softmax 2: The FFN Sigmoid (Inference)}

The second softmax is the sigmoid activation in the FFN. Its job is Bayesian
inference: what is my updated belief given my neighbors' beliefs?
\[
\updateBelief(m_0, m_1) = \sigma\!\left(\logit(m_0) + \logit(m_1)\right).
\]
This is where the actual Bayesian computation happens. The sigmoid is the exact
inverse of logit, converting the log-odds sum back to a probability. This is the
weight-of-evidence combination of Turing and Good~\cite{good1950probability},
applied to two binary evidence sources.

This softmax \emph{is} computing a posterior. The output is a belief in $(0,1)$
representing the updated probability that the current proposition is true, given
the gathered evidence. This is the inference step. Everything else in the
transformer is infrastructure for getting the right inputs to this computation.

\subsection*{Softmax 3: Output (Generation)}

The third softmax appears at the output layer, applied to the logits over the
vocabulary. Its job is generation: which token comes next? In a standard language
model, the logits are computed by a matrix multiply over the final hidden state ---
pattern-matched association rather than computed inference.

In a grounded QBBN transformer, the output would be different: the logits would come
from computed BP posteriors, representing the marginal probability of each
proposition being true given the evidence. The softmax operation is the same; what
changes is what it is applied to.

The QBBN does not replace the softmax. It replaces the logits. The architecture was
always capable of computing exact posteriors; it just needed the right weights to
feed into the output softmax. Standard language model training produces weights that
implement pattern-matched association. BP training would produce weights that
implement principled inference.

\subsection*{Why the Distinction Matters}

Conflating Softmax~1 and Softmax~2 leads to the mistaken conclusion that the
attention weights are the inference. They are not. The attention weights determine
\emph{which} evidence to gather. The FFN sigmoid determines \emph{what to conclude}
from that evidence. Routing and inference are separate computations, performed by
separate components, in a strict temporal order enforced by the architecture.

Conflating Softmax~2 and Softmax~3 leads to the mistaken conclusion that the output
distribution is a posterior over propositions. In an ungrounded model, it is not ---
it is a distribution over next tokens learned by association. The two converge only
in a grounded model where every output token corresponds to a declared proposition
with a known prior.
\section{OR Nodes vs.\ Dimensions: What $D_{ff}$ Is Actually Counting}
\label{sec:or-nodes-dimensions}

There is a persistent ambiguity in how to count the components of the implicit
factor graph. The formal construction has one OR gate per factor node, implemented
by one FFN hidden unit. A production transformer has $D_{ff} \approx 14{,}336$
hidden units per layer. Are there 14,336 factor nodes per layer? That seems too
many. This section resolves the ambiguity.

\subsection*{Hidden Units Are Not Factor Nodes}

In the formal construction, one factor node corresponds to one conditional
relationship between two propositions. One FFN hidden unit implements one such
relationship. So in the minimal construction, hidden units and factor nodes are
in one-to-one correspondence.

In a production transformer, this one-to-one correspondence breaks down in two
ways:

\begin{enumerate}
\item \textbf{Multiple units per factor.} A single conditional relationship may
require multiple hidden units to represent accurately with non-sigmoid activations.
ReLU and SwiGLU require more units to approximate the exact log-odds combination
that sigmoid implements in one unit. The factor node is still the unit of
semantics; the hidden units are the implementation.

\item \textbf{Superposed factors.} Because the residual stream represents many
concepts via superposition, a single hidden unit may participate in multiple
factor relationships simultaneously. The unit fires for a direction in residual
stream space that is a superposition of several feature directions. The
contribution it writes is a superposition of several belief updates.
\end{enumerate}

The right level of abstraction is not hidden units but \emph{SAE features}. Each
monosemantic SAE feature corresponds to one factor node in the implicit factor
graph. The number of factor nodes per layer is the number of active SAE features,
not $D_{ff}$.

\subsection*{What $D_{ff}$ Is Counting}

$D_{ff}$ is the \emph{implementation budget} for the OR gates, not the number of
factor nodes. It sets an upper bound on how many distinct conditional relationships
can be represented per layer, but the actual number of semantically distinct
relationships is determined by the SAE feature count, which is typically much
smaller than $D_{ff}$.

The ratio $D_{ff} / D_{\mathrm{model}} \approx 3$--$4$ reflects the implementation
overhead: roughly 3--4 hidden units are needed per residual stream dimension to
handle superposition, non-sigmoid activations, and $k$-ary neighbor chains. The
factor graph has far fewer nodes than $D_{ff}$ suggests; $D_{ff}$ is the engineering
budget for implementing those nodes accurately.

\subsection*{The Dimension Budget}

A cleaner way to think about the implicit factor graph uses the dimension budget
directly. The residual stream has $D_{\mathrm{model}}$ dimensions. Each active
concept occupies a direction in this space. The number of simultaneously active
concepts is bounded by the superposition capacity --- empirically, roughly
$10 \times D_{\mathrm{model}}$ features can be superposed with manageable
interference~\cite{elhage2022superposition}.

For Llama~3 405B with $D_{\mathrm{model}} = 16{,}384$, this gives roughly 160,000
simultaneously active features per token position. Each feature is a factor graph
node. The factor graph has on the order of $W \times 160{,}000 \approx 1.3\mathrm{B}$
nodes for a sequence of length 8192. This is the scale of the implicit Bayesian
network that one forward pass of Llama~3 405B is performing inference over.

\subsection*{The OR/AND Asymmetry Revisited}

With this clarification, the $D_{ff}/H \approx 416$ ratio has a cleaner
interpretation. It is not that there are 416 OR gates for every AND gate. It is
that the OR implementation requires 416 units of hidden dimension for every unit
of attention head dimension. The factor graph itself may be much sparser ---
each node may have only 2--5 neighbors in practice --- but implementing those
sparse connections accurately in the superposed representation requires the
$416\times$ overhead.

The implication for interpretability: finding the factor graph requires working at
the SAE feature level, not the hidden unit level. The factor potentials are encoded
in the relationship between input SAE features (what the $W_1$ rows respond to)
and output SAE features (what the $W_2$ columns write). The $D_{ff}$ hidden units
are the implementation substrate; the SAE features are the semantics.
\section{How to Read the Weights}
\label{sec:how-to-read-weights}

The previous sections have established what the transformer components
\emph{are} in Bayes net terms. This section asks a more practical question:
given a trained transformer, how would you actually find the factor graph?
What would you look for in the weight matrices?

The formal construction gives precise predictions. A transformer implementing
exact BP has weight matrices with specific sparse structure. A trained
transformer implementing approximate BP should show approximately this
structure. This section makes the predictions concrete enough to test.

\subsection*{What to Look for in Attention Weights}

The formal AND construction uses two weight families for each head:
\[
W_Q = W_K = \mathrm{projectDim}(d), \qquad W_V = \mathrm{crossProject}(s, d).
\]
$\mathrm{projectDim}(d)$ is a rank-1 matrix with a single nonzero entry on
the diagonal at position $(d,d)$. $\mathrm{crossProject}(s,d)$ is a rank-1
matrix with a single nonzero entry at position $(d,s)$.

The predictions for a semantic-AND head in a trained model:

\begin{enumerate}
\item \textbf{The Q/K circuit is approximately rank-1.} Compute the effective
Q/K matrix $W_Q^T W_K$ (or its low-rank approximation). A BP head should have
most of its singular value mass in a single component. The dominant left and
right singular vectors should correspond to the same interpretable SAE feature
--- the routing dimension.

\item \textbf{The dominant Q/K dimension is interpretable.} The routing
dimension $d$ should correspond to a monosemantic SAE feature: a direction in
residual stream space that activates for a single coherent concept. The head
attends to tokens that are high on this feature.

\item \textbf{The V circuit is approximately rank-1 and off-diagonal.} The
effective value matrix $W_V W_O$ should have most of its singular value mass
in a single component. The input direction (source dimension $s$) should
correspond to a belief-encoding feature (dimension~0 in the formal construction);
the output direction (destination dimension $d$) should correspond to a scratch
slot in the residual stream.

\item \textbf{The written dimension is distinct from the routing dimension.}
The head reads from dimension $d$ to determine whom to attend to, and writes
to dimension $d'$ to deposit the gathered belief. These should be different
SAE features --- routing and evidence are separate.
\end{enumerate}

\subsection*{What to Look for in FFN Weights}

The formal OR construction uses:
\[
W_1 \text{ row } i = \text{pattern that activates unit } i, \qquad
W_2 \text{ column } i = \text{belief update contributed by unit } i.
\]
With BP weights, the input pattern reads from the scratch slots (dimensions~4
and~5) and the output update writes to the belief dimension (dimension~0).

The predictions for FFN weights in a trained model:

\begin{enumerate}
\item \textbf{$W_1$ rows cluster by input feature.} Rows that read from the
same SAE feature should form a cluster. Each cluster corresponds to one factor
in the implicit factor graph.

\item \textbf{$W_2$ columns are interpretable belief updates.} Each column
should correspond to a human-interpretable concept update. ROME~\cite{meng2022rome}
already confirms this: individual columns encode factual associations.

\item \textbf{Input and output features are related by the factor graph.}
If $W_1$ row $i$ reads from features $A$ and $B$, then $W_2$ column $i$ should
update a feature $C$ where $C$ is a known consequence of $A$ and $B$ in the
knowledge domain. The factor $(A, B) \to C$ is the implicit conditional
probability table entry.
\end{enumerate}

\subsection*{The Closing Experiment}

The experiment that would close the loop:

\begin{enumerate}
\item Take a model with a trained SAE (e.g.\ Claude 3 Sonnet with the
Anthropic SAE).
\item Identify semantic-AND heads by behavior (sharp attention to a
content-determined neighbor, as in the psychic classification).
\item For each semantic-AND head, project $W_Q^T W_K$ onto the SAE feature
basis. Check if the dominant singular value component corresponds to a
monosemantic feature.
\item Check if the $W_V W_O$ circuit routes from a belief-encoding feature
to a scratch-slot feature.
\item For each active FFN unit, identify the input SAE features ($W_1$ row)
and the output SAE feature ($W_2$ column). Check if the relationship is
interpretable as a conditional probability table entry.
\end{enumerate}

If the results match the predictions, the formal BP construction has been
directly observed in a trained model. The implicit factor graph would be
legible: you could read off the factor potentials from the FFN weights and
the graph structure from the attention Q/K circuits.

\subsection*{Why This Has Not Been Done Yet}

The technical obstacle is SAE coverage. Current SAEs cover a small fraction
of the residual stream's representational capacity. The routing dimensions
predicted by the formal construction may not yet have clean SAE features
assigned to them. As SAE coverage improves --- and as the SAE feature basis
becomes more complete --- the experiment becomes more tractable.

The conceptual obstacle is that the field has not had a specific structural
prediction to test. The formal construction provides that prediction for the
first time. The weight inspection experiment is now a well-posed empirical
question, not a fishing expedition.
\section{The Hybrid BP Map: Four Head Types in Practice}
\label{sec:hybrid-bp}

The constructive BP mapping developed earlier applies to a specific idealized setting: sigmoid transformers with exact BP weights performing sharp index-matched routing. Production transformers differ from this in two important ways. They use activations such as GeLU or SwiGLU rather than sigmoid, and their attention heads exhibit a range of behaviors beyond sharp discrete lookup.

The right question is therefore not whether real transformers literally instantiate the constructive theorem in every component. It is whether the BP framework generates useful, testable interpretations of what trained transformers actually do. This section presents empirical head classification results and assesses what they support.

\subsection*{The Four Head Types}

Head classification across GPT-2 small, GPT-2 medium, and Qwen~2.5 0.5B, using unsupervised attention-pattern analysis over 88 prompts with no model-specific tuning, yields four stable functional categories:

\begin{center}
\begin{tabular}{llrrr}
\toprule
Type & Role & GPT-2 small & GPT-2 medium & Qwen 2.5 0.5B \\
\midrule
boolean-AND & Sharp routing / discrete gathering & 9 (6\%) & 9 (2\%) & 29 (9\%) \\
continuous-OR & Diffuse aggregation / soft evidence mixing & 28 (19\%) & 46 (12\%) & 59 (18\%) \\
hub & Global context communication & 103 (72\%) & 305 (79\%) & 203 (60\%) \\
mixed & Context-dependent behavior & 8 (6\%) & 24 (6\%) & 45 (13\%) \\
\bottomrule
\end{tabular}
\end{center}

\textbf{boolean-AND heads} attend to exactly one token with near-deterministic concentration. Their behavior matches the constructive routing heads of Section~\ref{sec:attention-and} most directly: one neighbor's information is gathered into a specific residual stream slot before the FFN update executes.

\textbf{continuous-OR heads} distribute attention broadly, integrating graded evidence from across the context window rather than performing discrete lookup. These heads are not well-described by the binary BP construction, but are naturally interpretable as soft evidence aggregators operating over continuous representations.

\textbf{Hub heads} consistently attend to or from position~0 (the BOS token) and constitute the dominant head population across all models studied. They are best understood as a special case of boolean-AND routing: the routing key always resolves to a fixed position rather than a content-determined neighbor. This makes them discrete routers like boolean-AND heads, just attending to a shared global workspace rather than a semantically matched token. The result is a star-shaped factor graph in which position~0 functions as a global variable node, accumulating evidence from all token positions and redistributing a global context signal back.

\textbf{Mixed heads} switch between sharp and diffuse behavior depending on input and do not exhibit a stable functional identity across prompts. They are concentrated in late layers in every model studied --- especially prominent in Qwen, where layers 22 and 23 collapse almost entirely to mixed behavior.

\subsection*{Coverage of the Framework}

Taken together, the three well-characterized types --- boolean-AND, continuous-OR, and hub --- account for approximately 90\% of all heads across every model studied (94\% in both GPT-2 models, 87\% in Qwen). Mixed heads, the genuinely uncharacterized remainder, constitute roughly 6--13\% depending on architecture. The BP-derived taxonomy covers the large majority of observed head behavior across model families.

\subsection*{Hub as Boolean-AND at a Fixed Address}

Because hub heads are a special case of boolean-AND routing, the AND and hub populations together form the discrete routing infrastructure of the network. Combining them:

\begin{center}
\begin{tabular}{lrrr}
\toprule
 & GPT-2 small & GPT-2 medium & Qwen 2.5 0.5B \\
\midrule
discrete routing (AND + hub) & 112 (78\%) & 314 (82\%) & 232 (69\%) \\
continuous aggregation (OR) & 28 (19\%) & 46 (12\%) & 59 (18\%) \\
uncharacterized (mixed) & 8 (6\%) & 24 (6\%) & 45 (13\%) \\
\bottomrule
\end{tabular}
\end{center}

Discrete routing dominates in every model. The architecture is primarily a routing system, with a minority of heads performing continuous evidence aggregation.

\subsection*{The Scaling of boolean-AND Heads}

Within the GPT-2 family, boolean-AND head count is stable: GPT-2 small (124M parameters, 144 total heads) contains 9, and GPT-2 medium (345M parameters, 384 total heads) also contains exactly 9. Total head count grows 2.8$\times$; sharp routing infrastructure does not. What grows instead is hub infrastructure, from 72\% to 79\% of all heads.

Across architectures, the count varies: Qwen 2.5 0.5B at a comparable scale has 29 boolean-AND heads. This suggests that the number of discrete routing heads reflects architectural choices rather than a universal constant. Different model families allocate different amounts of capacity to sharp content-based routing. The consistent finding across all three models is that boolean-AND heads remain a small minority --- the routing primitive is compact regardless of how many instances a given architecture deploys.

This is consistent with the BP interpretation: routing is a compact operation. What the architecture scales is inference capacity --- hub infrastructure and OR aggregation --- not the routing primitive itself.

\subsection*{Hybrid Inference}

The constructive theorem establishes exact BP realization for sigmoid transformers. Production transformers use SwiGLU or GeLU. Does the BP interpretation break down?

Not entirely --- but the character of the inference changes. Bayesian networks are not limited to binary variables. Continuous-variable Bayesian networks, Gaussian belief propagation, and expectation propagation are all well-established frameworks in which message passing operates over graded rather than discrete representations.

The empirical head distribution suggests a hybrid picture: boolean-AND and hub heads implement discrete routing within the BP family; continuous-OR heads implement inference-like operations over graded representations that the binary construction does not cover exactly. The constructive sigmoid theorem is then best understood as the exact limiting case of a broader family of message-passing computations that production transformers approximate.

For the boolean-AND heads specifically, the formal predictions are testable: their Q/K circuits should exhibit approximately rank-1 structure consistent with \texttt{projectDim}, and their V circuits should show \texttt{crossProject}-style routing. This is the experiment that would most directly confirm the BP interpretation in trained models.

\subsection*{Layer Structure}

Head type distribution is not uniform across layers. GPT-2 small exhibits a three-phase structure that aligns naturally with iterative inference:

\begin{itemize}
\item \textbf{Layers 0--2} are dominated by continuous-OR heads performing local feature extraction. Boolean-AND heads are present but sparse. This corresponds to early rounds of BP: each node gathers evidence from nearby tokens.

\item \textbf{Layer 3} is a transition layer in which all four head types are represented. Boolean-AND heads are active at exactly the point where the network shifts toward global communication.

\item \textbf{Layers 4--11} are hub-dominated. The primary computation shifts to reading from and writing to the global context at position~0.
\end{itemize}

This progression --- local evidence gathering followed by global integration --- is what you would expect from an iterative inference system. It is consistent with the BP interpretation but not uniquely predicted by it; other accounts of transformer computation could produce a similar layer structure.

\subsection*{Open Questions}

Three questions mark the boundary of what the hybrid BP account currently explains:

\begin{enumerate}
\item \textbf{The continuous update rule.} What is the correct analog of $\updateBelief$ for continuous nodes? GeLU and SwiGLU preserve non-negativity and smooth aggregation, but their relationship to Gaussian BP or expectation propagation has not been worked out.

\item \textbf{What determines the AND head count?} Boolean-AND head count is stable within the GPT-2 family but varies across architectures --- 9 in GPT-2, 29 in Qwen at comparable scale. What architectural choices drive this, and what is the minimum discrete routing infrastructure required for language modeling?

\item \textbf{Prediction vs.\ compatibility.} The BP framework is consistent with many observed transformer behaviors. The stronger claim --- that it \emph{predicts} them --- requires identifying behaviors the framework would rule out if false. Specifying those tests is the next step toward treating the hybrid BP account as a genuine scientific hypothesis rather than an interpretive lens.
\end{enumerate}

The constructive theorem gives an exact reference point for the boolean-AND case. The continuous and mixed populations remain the empirical frontier.

\part{Grounding and Hallucination}
\section{Grounded vs.\ Ungrounded Transformers}
\label{sec:grounded-ungrounded}

A grounded transformer and a standard language model may share the same
underlying architecture while differing substantially in how their
internal representations relate to externally defined concepts and
verification procedures.

Within the BP interpretation developed in this paper, both grounded and
ungrounded systems perform inference over implicit factor-graph-like
structures. The key distinction is whether the graph has a formally
declared interpretation.

\subsection*{The Distinction}

A \emph{grounded} transformer has:

\begin{itemize}
\item An explicit factor graph: tokens or latent states correspond to
declared propositions in a known knowledge base.

\item A finite concept space: the relevant propositions are bounded and
specified in advance.

\item A verifier: outputs can be checked against externally defined
correctness criteria.

\item Formal guarantees in restricted settings: on tree-structured
knowledge bases with constructive BP weights, exact posterior guarantees
follow from the companion theorem~\cite{coppola2026transformers}.
\end{itemize}

An \emph{ungrounded} transformer has:

\begin{itemize}
\item An implicit factor graph encoded in learned parameters rather than
explicitly declared structures.

\item No formally declared ontology tying internal states to a finite,
verified concept inventory.

\item No architecture-internal correctness guarantee analogous to the
grounded setting.

\item Reliability that is primarily empirical rather than formally
verified.
\end{itemize}

\subsection*{Hallucination in BP Terms}

Within this framework, hallucination can be interpreted as a mismatch
between model beliefs and externally verified posteriors.

For grounded systems, the comparison is well-defined:
\[
b(j) \neq P_{\mathrm{true}}(j),
\]
where $P_{\mathrm{true}}(j)$ is the correct posterior induced by the
knowledge base and observed evidence.

Ungrounded language models differ because they lack a formally declared
ontology and verifier. Their outputs may still correspond to meaningful
or true propositions in ordinary semantic practice, but the architecture
does not itself provide a mechanism guaranteeing correctness relative to
a fixed external knowledge base.

From this perspective, hallucination is not merely a numerical error in
the inference procedure. It is also a consequence of operating without a
fully specified and verifiable grounding structure.

\subsection*{Scaling and Grounding}

Scaling improves empirical performance. Larger models generally produce
more coherent outputs, achieve lower hallucination rates on benchmarks,
and exhibit richer internal representations.

However, scaling alone does not provide the same kind of formal
correctness guarantees available in grounded settings. Increasing model
capacity improves the quality of the learned implicit factor graph, but
it does not by itself create an explicit ontology, verifier, or formally
grounded concept space.

The distinction is therefore not between ``working'' and ``non-working''
systems. It is between empirical reliability and formally specified
correctness guarantees.

\subsection*{What Grounding Would Look Like}

A fully grounded transformer system would require several components:

\begin{enumerate}
\item A declared concept space: propositions or entities are explicitly
represented and bounded.

\item A declared relational structure: dependencies between concepts are
specified or verified.

\item A grounding procedure: natural language inputs are mapped into the
declared representational system.

\item A verification mechanism: outputs can be checked against the
underlying ontology or knowledge base.
\end{enumerate}

Retrieval-augmented generation (RAG) can be viewed as a partial form of
grounding because retrieved documents temporarily constrain the inference
space. More complete grounding would require persistent ontological and
verification structure rather than transient context retrieval alone.

\subsection*{The Practical Implication}

The practical distinction between grounded and ungrounded systems is the
kind of reliability guarantee they support.

For grounded systems in restricted settings, the constructive BP results
yield exact posterior guarantees. For general-purpose language models,
the strongest guarantees are typically empirical: benchmark performance,
human evaluation, calibration studies, or retrieval consistency.

The gap between these two forms of reliability --- formal verification
versus empirical robustness --- is the grounding gap.

\part{Convergent Evidence}
\section{Evidence from Mechanistic Interpretability}
\label{sec:interp-evidence}

The mechanistic interpretability literature provides some of the
strongest empirical motivation for viewing transformers through a
message-passing lens. Importantly, most of this work was not originally
framed in terms of belief propagation or Bayesian networks. The value of
the BP perspective is therefore not that it replaces the
interpretability literature, but that it offers a common language for
organizing many of its findings.

This section surveys several influential interpretability results and
discusses how they naturally fit into a routing-and-update view of
transformer computation.

\subsection*{Induction Heads}

Olsson et al.~\cite{olsson2022induction} identify ``induction heads'':
two-head circuits that implement in-context pattern matching. Head~1
(the ``prefix'' head) attends from position $t$ back to the previous
occurrence of the current token. Head~2 (the ``induction'' head) then
attends to the token that followed that previous occurrence.

From the tutorial BP perspective, this resembles a two-stage gather
operation. One head retrieves a historical context; the second uses that
retrieved context to identify the next relevant source of evidence. The
important point is not that induction heads are ``literally'' BP nodes,
but that the routing structure closely resembles iterative message
passing over an implicit dependency graph.

\subsection*{Name Mover Heads and the IOI Circuit}

Wang et al.~\cite{wang2022ioi} reverse-engineer the indirect object
identification circuit in GPT-2 small. The key components are ``name
mover heads'' that copy token identity information from one position to
another. The Q/K circuits of these heads exhibit approximately rank-1
structure, while the V circuits route token-specific content into target
positions.

This behavior is naturally compatible with the routing interpretation
developed earlier in the paper. In the constructive BP mapping, sparse
projectDim/crossProject-style operations perform a similar function:
matching one feature dimension and routing information into another
location in the residual stream.

The BP framework therefore provides one possible interpretation of why
such sparse routing patterns repeatedly emerge in trained models.

\subsection*{FFN Key-Value Memories}

Geva et al.~\cite{geva2021ffn} show that FFN layers behave like
key-value memories: the first-layer weight matrix detects patterns and
the second-layer matrix writes structured updates back into the residual
stream. Many hidden units correspond to human-interpretable concepts or
contexts.

Within the tutorial framework developed here, FFNs can be interpreted as
performing local update operations conditioned on gathered evidence. The
``key'' determines when a particular update should activate; the
``value'' determines what modification is written into the evolving
residual-stream state.

This interpretation does not require treating FFN rows as literal
conditional-probability tables. Rather, it suggests that the key-value
memory view and the message-passing view may be complementary ways of
describing the same computational structure.

\subsection*{ROME and Localized Editing}

Meng et al.~\cite{meng2022rome} demonstrate that factual associations in
GPT models can sometimes be edited by modifying localized FFN
parameters. The resulting edits are often surprisingly surgical:
changing one association without broadly disrupting neighboring
behavior.

From the BP perspective, this resembles modifying a localized update
rule or factor potential within a larger inference system. Again, the
point is not that ROME directly proves the existence of explicit Bayes
net tables inside the transformer. The point is that localized,
structured update behavior is naturally compatible with a factorized
message-passing interpretation.

\subsection*{Sparse Autoencoders and Features}

Templeton et al.~\cite{templeton2024monosemanticity} train sparse
autoencoders on Claude 3 Sonnet and recover millions of interpretable
features from the residual stream. Many features exhibit coherent
semantic structure, feature neighborhoods, and suppression
relationships.

Within the tutorial framing, these features can be viewed as candidate
nodes or directions in an implicit computational graph. Their activation
values behave like evolving internal state variables whose interactions
are mediated through attention and FFN updates.

The ``Golden Gate Claude'' experiment is especially suggestive from this
perspective: clamping one feature propagates structured downstream
effects through the model, much like injecting evidence into a connected
message-passing system.

\subsection*{Attribution Graphs}

Recent work from Anthropic~\cite{anthropic2025attribution} constructs
attribution graphs tracing information flow between features across
layers. These graphs make visible the pathways through which activations
influence one another during a computation.

From the tutorial BP perspective, attribution graphs resemble partial
visualizations of the model's implicit message-passing structure:
attention routes information, FFNs apply transformations, and residual
streams accumulate evolving state over depth.

\subsection*{Synthesis}

The BP perspective should be understood as an organizing framework for
these findings rather than a claim that interpretability researchers
were unknowingly ``discovering Bayes nets.''

The important observation is simply that many empirical findings in
mechanistic interpretability fit naturally into a routing-plus-update
view of transformer computation:
\begin{itemize}
\item induction heads resemble multi-stage retrieval,
\item routing heads move structured information between positions,
\item FFNs apply localized updates,
\item sparse features behave like persistent latent variables,
\item and attribution graphs reveal structured pathways of influence.
\end{itemize}

The convergence between these findings and the message-passing
perspective is one reason the BP framework is useful as a conceptual
tool for understanding transformers.
\section{Related Work}
\label{sec:related-work}

Five independent research programs have converged on the conclusion that the
transformer forward pass implements belief propagation. None of them set out to
prove this. Each approached the transformer from a different direction and found
the same structure. This section surveys each program and its relationship to
the account developed in this paper.

\subsection*{Belief State Geometry}

Several recent papers establish that transformer residual streams represent belief
states linearly. Shai et al.~\cite{shai2024belief} show that the residual stream
encodes a linear representation of the agent's belief state over possible world
states, updated by attention operations. Piotrowski et
al.~\cite{piotrowski2025constrained} show that attention implements constrained
belief updates consistent with Bayesian conditioning. Agarwal et
al.~\cite{agarwal2025bayesian} show that transformers are primitively complete for
Bayesian inference --- any Bayesian computation can be implemented by a transformer.

These results are complementary to the account here. They establish the
\emph{what} (the residual stream is a belief state; attention implements Bayesian
updates) without fully specifying the \emph{how} (the specific weight construction
that makes this exact, the AND/OR decomposition, the factor graph structure). The
formal construction of~\cite{coppola2026transformers} provides the missing
mechanistic account.

\subsection*{In-Context Learning as Bayesian Inference}

Xie et al.~\cite{xie2022icl} argue that in-context learning is implicit Bayesian
inference at the population level: the model maintains a posterior over latent
document concepts and updates it as it sees more examples. This is a statistical
account --- it describes what ICL accomplishes in terms of the training distribution
--- rather than a mechanistic account of how individual forward passes implement it.

The BP account is the mechanistic complement: each forward pass is one round of BP
on the implicit factor graph, and in-context learning works because the implicit
factor graph encodes the conditional relationships learned from training. The two
accounts are consistent and describe different levels of abstraction.

\subsection*{Constructive Weight Proofs}

Aky\"{u}rek et al.~\cite{akyurek2023icl} and von Oswald et
al.~\cite{oswald2023gradient} show that specific weight patterns make transformers
perform implicit gradient descent, explaining in-context learning as approximate
optimization. The methodology is the same as~\cite{coppola2026transformers}:
construct explicit weights that implement a specific algorithm, prove correctness,
confirm empirically.

The key differences: BP is exact on trees (not approximate or asymptotic), and the
BP account is mechanistic per-instance (each forward pass computes a specific
posterior) rather than statistical at the population level (the model behaves as if
it were doing Bayesian inference in expectation over the training distribution).

\subsection*{GNNs, Message Passing, and Attention}

Scarselli et al.~\cite{scarselli2009gnn} established that graph neural networks
implement belief propagation on explicit graph structure. Gilmer et
al.~\cite{gilmer2017mpnn} unified GNNs under the message passing neural network
(MPNN) framework. Yoon et al.~\cite{yoon2019gnn} extended this to approximate
inference. Joshi~\cite{joshi2020gnn} observed that full self-attention is equivalent
to a GNN on a complete graph.

The contribution of~\cite{coppola2026transformers} is orthogonal: a standard
transformer with full self-attention and no attention mask implements BP on an
explicit \emph{external} factor graph, where graph structure is encoded in token
embeddings and routing is discovered via Q/K index matching. The graph is not
hardwired into the attention mask --- it is represented in the token content.

\subsection*{Mechanistic Interpretability}

Elhage et al.~\cite{elhage2021mathematical} develop the mathematical framework for
transformer circuits. Olsson et al.~\cite{olsson2022induction} identify induction
heads as a specific named algorithm discovered by gradient descent. Wang et
al.~\cite{wang2022ioi} reverse-engineer the full IOI circuit in GPT-2 small.
Geva et al.~\cite{geva2021ffn} establish the key-value memory interpretation of
FFN layers. Meng et al.~\cite{meng2022rome} show that factual associations can be
edited by modifying FFN weights. Templeton et al.~\cite{templeton2024monosemanticity}
find millions of monosemantic features via SAEs.

The BP account gives this literature a unified theoretical framework. Every major
finding maps onto a component of the belief propagation account: induction heads
are two-round BP gathers; name mover heads are projectDim/crossProject routing; FFN
key-value memories are factor potentials; ROME edits are factor potential
modifications; SAE features are factor graph nodes. The field has been mapping the
implicit factor graph without a map. This paper provides the map.

\subsection*{The Log-Odds Tradition}

Good~\cite{good1950probability} and Turing formalized log-odds addition as the
correct algebra for combining independent binary evidence. Pearl~\cite{pearl1988probabilistic}
applied this algebra to graphical models via the sum-product algorithm. The
connection between this tradition and the transformer architecture has not previously
been made explicit before~\cite{coppola2026transformers}. Section~\ref{sec:log-odds}
of this paper traces the lineage and establishes that the sigmoid transformer is the
mechanical implementation of the Turing-Good-Pearl algebra.

\subsection*{Summary}

The BP account of the transformer is not a new claim in isolation. It is the
convergence point of five independent research programs. The formal
contribution of~\cite{coppola2026transformers} is to make the convergence
rigorous: a formally verified constructive proof that a sigmoid transformer with
any weights implements weighted BP on its implicit factor graph, with explicit
weights for the exact case, a uniqueness theorem showing exact posteriors force
those weights, and empirical confirmation that gradient descent finds the
predicted structure. This paper makes that formal account concrete and readable.

\bibliographystyle{plain}
\bibliography{refs}

\end{document}